\chapter{Extended Related Work Analysis}
\label{appendix:related_works}

\section{Single-Tenant Agent Frameworks}

\subsection{MemGPT (Packer et al., 2023)}

\textbf{Core Contribution}: Hierarchical memory system with OS-like virtual context management.

\textbf{Architecture}:
\begin{itemize}
    \item \textbf{Main Context}: Active LLM working memory
    \item \textbf{Disk}: Long-term storage for conversation history
    \item \textbf{Paging}: Dynamically swap memory pages between main context and disk
\end{itemize}

\textbf{Limitations for Networked Scenarios}:
\begin{itemize}
    \item Memory is globally accessible—no relationship-specific partitioning
    \item No security boundaries—assumes trusted single-tenant execution
    \item No cross-agent communication primitives
\end{itemize}

\textbf{Pulse's Advancement}: Dual-Track Memory extends MemGPT's tiered architecture with relationship shards, enabling secure multi-party interactions.

\subsection{OpenClaw / Manus (Commercial Systems)}

\textbf{Core Contribution}: Production-grade task automation with computer vision and browser control.

\textbf{Capabilities}:
\begin{itemize}
    \item GUI automation (clicking, typing, navigating)
    \item Multi-step workflow orchestration
    \item Integration with desktop applications
\end{itemize}

\textbf{Limitations for Networked Scenarios}:
\begin{itemize}
    \item Designed for internal workspace—no access governance model
    \item Assumes user is always the principal—no external agent concept
    \item Security through obscurity (user's local machine)
\end{itemize}

\textbf{Pulse's Advancement}: Context Cells provide physical security boundaries enabling external parties to interact with the agent without full workspace access.

\subsection{AutoGPT / BabyAGI (Open Source)}

\textbf{Core Contribution}: Autonomous task decomposition and execution loops.

\textbf{Architecture}:
\begin{itemize}
    \item Self-prompting agents that plan multi-step tasks
    \item Tool use through plugins
    \item Memory through vector databases or simple files
\end{itemize}

\textbf{Limitations for Networked Scenarios}:
\begin{itemize}
    \item No persistent identity—each run starts fresh
    \item No relationship awareness
    \item No security model for external access
\end{itemize}

\textbf{Pulse's Advancement}: Proactive Heartbeat mechanism + persistent identity + relationship-aware memory create a true "lifelong" agent, not just a task executor.

\section{Agent Communication Protocols}

\subsection{Multi-Agent Systems (JADE, FIPA-ACL)}

\textbf{Core Contribution}: Standardized message-passing protocols for agent coordination.

\textbf{FIPA-ACL Message Structure}:
\begin{verbatim}
(REQUEST
  :sender agent1
  :receiver agent2
  :content (action (deliver package123))
  :language FIPA-SL)
\end{verbatim}

\textbf{Design Philosophy}:
\begin{itemize}
    \item Agents as autonomous software entities
    \item Speech acts for communication (REQUEST, INFORM, QUERY)
    \item Ontology-based content representation
\end{itemize}

\textbf{Limitations for LLM-Based Agents}:
\begin{itemize}
    \item Assumes deterministic agent behavior—no handling of LLM non-determinism
    \item No security model—assumes trusted agent federation
    \item Rigid ontologies—incompatible with natural language flexibility
\end{itemize}

\textbf{Pulse's Advancement}: Hybrid approach—structured JSON payloads for reliability + natural language for flexibility, combined with zero-trust authentication via Context Cells.

\subsection{LLM-Based A2A Communication (Wang et al., 2024)}

\textbf{Core Contribution}: Agents negotiate through natural language exchanges.

\textbf{Example Flow}:
\begin{verbatim}
Agent A: "I propose we meet on Tuesday at 2pm."
Agent B: "I'm unavailable then. How about Wednesday at 3pm?"
Agent A: "Wednesday works. Confirmed."
\end{verbatim}

\textbf{Advantages}:
\begin{itemize}
    \item Leverages LLM's natural language understanding
    \item Flexible—can handle unanticipated scenarios
\end{itemize}

\textbf{Limitations}:
\begin{itemize}
    \item High latency (3+ LLM roundtrips per agreement)
    \item Token-inefficient (verbose text exchanges)
    \item Hallucination risks (Agent A might agree to something Agent B never proposed)
    \item No access governance—assumes agents have unrestricted context
\end{itemize}

\textbf{Pulse's Advancement}: Structured intent exchange for common tasks (scheduling, document requests) with natural language fallback for complex negotiations. Context Cells ensure data access is scoped per relationship.

\section{Privacy-Preserving AI Systems}

\subsection{Federated Learning}

\textbf{Core Contribution}: Train models collaboratively without centralizing data.

\textbf{Architecture}:
\begin{itemize}
    \item Local model training on device
    \item Gradient sharing (not raw data)
    \item Central aggregation for global model
\end{itemize}

\textbf{Relevance to Pulse}:
\begin{itemize}
    \item Similar philosophy: enable collaboration without full data sharing
    \item Different use case: FL is for model training, Pulse is for operational coordination
\end{itemize}

\textbf{Potential Integration}: Federated learning could enable agents to improve their policy models (Escalation Protocol) without sharing sensitive precedent data with a central server.

\subsection{Differential Privacy}

\textbf{Core Contribution}: Quantify and limit information leakage from queries.

\textbf{Definition}:
\[
P[M(D) \in S] \leq e^{\epsilon} \cdot P[M(D') \in S] + \delta
\]

where $D$ and $D'$ differ by one record.

\textbf{Application to Pulse}: Future work could formally bound the information leakage from indirect inference attacks using differential privacy metrics.

\textbf{Challenge}: DP typically adds noise to results. For operational coordination (e.g., ``When are you free?''), noisy answers are unacceptable. Research needed on \textit{selective DP}—apply DP only to sensitive aggregates, not individual operational queries.

\section{AI Safety and Alignment}

\subsection{Constitutional AI (Anthropic)}

\textbf{Core Contribution}: Self-critique and revision based on constitutional principles.

\textbf{Training Process}:
\begin{enumerate}
    \item Generate initial response
    \item Critique response against constitution (``Is this helpful and harmless?'')
    \item Revise response based on critique
    \item Train model on revised responses (RLHF)
\end{enumerate}

\textbf{Relevance to Pulse}:
\begin{itemize}
    \item The ``Macro Constitution'' in Pulse's Escalation Protocol is inspired by Constitutional AI
    \item However, Pulse enforces constitution through \textit{system architecture}, not just training
\end{itemize}

\textbf{Key Difference}: Constitutional AI trusts the model to follow principles. Pulse's Context Cells provide physical enforcement—the model \textit{cannot} violate the constitution even if compromised.

\subsection{RLHF and Instruction Following}

\textbf{Core Contribution}: Fine-tune models to follow user instructions reliably.

\textbf{Limitations Exposed by Pulse}:
\begin{itemize}
    \item RLHF improves average-case instruction following but does not eliminate adversarial brittleness
    \item Jailbreak research (Zou et al., 2023; Wei et al., 2024) demonstrates that RLHF-aligned models remain vulnerable
    \item Pulse's approach: \textit{assume the model is adversarial} and design the system accordingly
\end{itemize}

\textbf{Future Integration}: RLHF could improve Pulse's Sanitization Agent (fewer false positives in Escalation Protocol), but security guarantees must remain architectural, not behavioral.

\section{Data Sharing and Collaboration Platforms}

\subsection{Notion, Confluence, Google Workspace}

\textbf{Core Contribution}: Collaborative document editing with permission management.

\textbf{Permission Models}:
\begin{itemize}
    \item \textbf{File-Level ACLs}: Grant read/write access per document
    \item \textbf{Sharing Links}: Time-limited, password-protected URLs
    \item \textbf{Team Workspaces}: Organization-wide visibility with role-based access
\end{itemize}

\textbf{Limitations}:
\begin{itemize}
    \item Static—documents don't answer questions or execute actions
    \item Manual—users must explicitly share each resource
    \item No intelligence—no understanding of relationship context or precedents
\end{itemize}

\textbf{Pulse's Advancement}: Context Cells bring file ACLs to the agent execution layer. Sharing a Pulse "link" grants access not just to static files but to an intelligent agent scoped to those files.

\subsection{Dropbox, Google Drive}

\textbf{Core Contribution}: Secure file storage and sharing at scale.

\textbf{Security Features}:
\begin{itemize}
    \item Encryption at rest and in transit
    \item Granular permissions (view/comment/edit)
    \item Audit logs for access tracking
\end{itemize}

\textbf{Pulse's Conceptual Parallel}:
\begin{itemize}
    \item Dropbox shares \textit{files}
    \item Pulse shares \textit{agentic access} to files
\end{itemize}

\textbf{Future Vision}: Pulse as a "Dropbox for Agents"—a file system where access is mediated not by direct downloads but by AI proxies operating within Context Cells.

\section{Conversational AI and Personalization}

\subsection{Replika, Character.AI}

\textbf{Core Contribution}: Long-term conversational companions with personality and memory.

\textbf{Personalization Mechanisms}:
\begin{itemize}
    \item \textbf{Profile-Based}: User explicitly defines persona (gender, interests, communication style)
    \item \textbf{Memory-Based}: Agent recalls past conversations to maintain continuity
    \item \textbf{Emotion Modeling}: Agents respond to user emotional state
\end{itemize}

\textbf{Limitations}:
\begin{itemize}
    \item Synthetic relationships—no grounding in real-world operational state
    \item No tool use—agents cannot execute actions
    \item Single-tenant—designed for 1:1 user-agent interaction, no multi-party scenarios
\end{itemize}

\textbf{Pulse's Advancement}: Combines the personalization depth of Replika (Relationship Shards) with real-world grounding (calendar, email, files) and multi-party security (Context Cells).

\section{Enterprise AI Assistants}

\subsection{Microsoft Copilot, Google Duet AI}

\textbf{Core Contribution}: AI assistance integrated into productivity suites.

\textbf{Capabilities}:
\begin{itemize}
    \item Summarize emails, draft responses
    \item Generate presentation slides from documents
    \item Analyze spreadsheets and suggest insights
\end{itemize}

\textbf{Security Model}:
\begin{itemize}
    \item Tenant isolation—each organization's data is firewalled
    \item Role-based access—Copilot respects existing permissions (e.g., cannot read emails user doesn't have access to)
\end{itemize}

\textbf{Limitations for Cross-Org Coordination}:
\begin{itemize}
    \item Siloed—Microsoft Copilot for Company A cannot interact with Google Duet AI for Company B
    \item No federated identity—cross-organizational collaboration still requires manual email/calendar coordination
\end{itemize}

\textbf{Pulse's Advancement}: Network-native design enables cross-organizational agent interactions with cryptographic delegation protocols (Transitive Delegation).

\section{Theoretical Frameworks}

\subsection{Capability-Based Security (Dennis and Van Horn, 1966)}

\textbf{Core Principle}: Objects are accessible only through unforgeable tokens (capabilities).

\textbf{Classic Example}: UNIX file descriptors
\begin{itemize}
    \item Process receives file descriptor (integer token) after successful \texttt{open()}
    \item Descriptor grants access—no need to re-authenticate on each \texttt{read()}
    \item Kernel enforces that processes cannot forge or steal descriptors
\end{itemize}

\textbf{Pulse's Application}:
\begin{itemize}
    \item Context Cells are capabilities for agent execution
    \item External guests receive \texttt{agenticLink} tokens that encode allowed resources
    \item PDP enforces capability validity at runtime—agents cannot escalate privileges
\end{itemize}

\subsection{Information Flow Control (Denning, 1976)}

\textbf{Core Principle}: Track data provenance to prevent unauthorized leakage.

\textbf{Security Lattice}:
\[
L = (S, \leq, \sqcup, \sqcap)
\]

where $S$ is a set of security labels, $\leq$ is the dominance relation, and $\sqcup$, $\sqcap$ are join/meet operations.

\textbf{Application to Pulse}:
\begin{itemize}
    \item Each file/state has a security label (e.g., \texttt{confidential\_financials})
    \item Tool outputs inherit labels from inputs (taint propagation)
    \item Responses to external guests are checked against allowed labels before sending
\end{itemize}

\textbf{Future Work}: Implement formal IFC tracking to detect multi-query inference attacks (Section~\ref{chapter:evaluation}).
